%%%%%%%%%%%%%%%%%%%%%%%%%%%%%%%%%%%%%%%%%%%%%%%%%%%%%%%%%%%%
\section*{Appendix C: Framework Integration Details}
\label{appendix:framework-details}

{\scriptsize
Detailed mappings showing how framework constructs map to authenticated workflow primitives across nine frameworks.

\subsection*{A2A (Agent-to-Agent Protocol)}

A2A defines standards for agent discovery, trust establishment, and communication. MACAW provides the cryptographic infrastructure A2A requires:

{\tiny
\begin{verbatim}
A2A Requirements:              MACAW Already Provides:
┌──────────────────┐           ┌───────────────────────┐
│Agent Discovery   │  ────►    │Agent Registry         │
│                  │           │- Agent directory      │
│                  │           │- Capability ads       │
├──────────────────┤           ├───────────────────────┤
│Identity Verify   │  ────►    │MA (Crypto Identity)   │
│                  │           │- Public/private keys  │
│                  │           │- X.509 certificates   │
├──────────────────┤           ├───────────────────────┤
│Trust Establish   │  ────►    │Attestation System     │
│                  │           │- Claims about agent   │
│                  │           │- Verifiable proofs    │
├──────────────────┤           ├───────────────────────┤
│Secure Messaging  │  ────►    │CAW (Signed Invocation)│
│                  │           │- Message signing      │
│                  │           │- Replay protection    │
├──────────────────┤           ├───────────────────────┤
│Protocol Negotiat │  ────►    │Capability Exchange    │
│                  │           │- Version negotiation  │
│                  │           │- Feature discovery    │
└──────────────────┘           └───────────────────────┘

Complete A2A Flow:
┌───────────┐                          ┌───────────┐
│ Agent A   │  1. Discovery Request    │ Agent B   │
│           │ ────────────────────────►│           │
│           │  "Who can analyze data?" │           │
└─────┬─────┘                          └─────┬─────┘
      │                                      │
      ▼                                      ▼
┌──────────────────────────────────────────────────┐
│              Agent Registry                      │
│  Registry:                                       │
│   - Agent B: capabilities["data_analysis"]       │
│   - Agent B: public_key="..."                    │
│   - Agent B: attestations["certified_analyst"]   │
└──────────────────────────────────────────────────┘
      │ 2. Returns Agent B info
      ▼
┌───────────┐  3. Establish trust      ┌───────────┐
│ Agent A   │ ────────────────────────►│ Agent B   │
│           │  - Verify B's signature  │           │
│           │  - Check attestations    │           │
│           │  - Validate policy       │           │
│           │                          │           │
│           │  4. Send CAW             │           │
│           │ ────────────────────────►│           │
│           │  - Message signed by A   │           │
│           │  - A's attestations      │           │
│           │  - Policy requirements   │           │
└───────────┘                          └───────────┘
\end{verbatim}}

\textbf{Security Challenge}: Agents must verify identity and capabilities before delegating authority. Attestations provide cryptographic claims that enable trust decisions without centralized authorities.

\subsection*{MCP (Model Context Protocol)}

MCP servers expose three capability types: tools (executable functions), resources (data access with read/write/watch operations), and prompts (retrievable templates). The MACAW mapping:

{\tiny
\begin{verbatim}
MCP Core Concepts:
┌─────────────────────────────────────────────┐
│              MCP Server                     │
├──────────┬──────────────┬───────────────────┤
│  Tools   │  Resources   │     Prompts       │
│execute() │ read/write() │ getPrompt()       │
└──────────┴──────────────┴───────────────────┘

MACAW Mapping:
┌─────────────────────────────────────────────┐
│         macawAgent (MCP Server)             │
├──────────┬──────────────┬───────────────────┤
│  Tool    │ Tool         │ Authenticated     │
│  (execute│ (Resource)   │ Prompt            │
│   → CAW) │ (read → CAW) │ (signed template) │
└──────────┴──────────────┴───────────────────┘

Resource Flow:
Client: "Read file.txt" → SecureMCP: invoke_tool()
  → PEP: signature valid? → Policy: allow read?
  → Execute if allowed
\end{verbatim}}

\textbf{Security Challenge}: Resources represent read-only data access, while tools can modify state. The mapping enforces read-only constraints through policy, preventing privilege escalation from data retrieval to data modification.

\subsection*{OpenAI and Claude (LLM Interfaces)}

LLM providers expose inference APIs with function calling. Two deployment variants exist based on trust assumptions:

{\tiny
\begin{verbatim}
Two-sided Wrapping: Both LLM API and Client as macawAgents
┌─────────────────────────────────────────────────┐
│            Application Code                     │
└───────────────┬─────────────────────────────────┘
                ▼
┌─────────────────────────────────────────────────┐
│          SecureOpenAI/SecureClaude              │
├────────────────────┬────────────────────────────┤
│  LLM Wrapper       │   Tool Wrapper             │
│                    │                            │
│ chat.complete()    │  search_tool()             │
│      ↓             │       ↓                    │
│  INTERCEPT         │   INTERCEPT                │
│      ↓             │       ↓                    │
│  Sign as macawAgent│   Sign invocation          │
│      ↓             │       ↓                    │
│  PEP verifies:     │   PEP verifies:            │
│  - Signature       │   - Signature              │
│  - Provider policy │   - App policy             │
│      ↓             │       ↓                    │
│  Execute LLM call  │   Execute tool             │
└────────────────────┴────────────────────────────┘

One-sided Wrapping: Only Client Tools as macawAgents
┌─────────────────────────────────────────────────┐
│            Application Code                     │
└───────────────┬─────────────────────────────────┘
                ▼
┌─────────────────────────────────────────────────┐
│          SecureOpenAI/SecureClaude              │
├────────────────────┬────────────────────────────┤
│  LLM Passthrough   │   Tool Wrapper             │
│                    │                            │
│ chat.complete()    │  search_tool()             │
│      ↓             │       ↓                    │
│  Direct to API     │   INTERCEPT                │
│  (no interception) │       ↓                    │
│                    │   Sign invocation          │
│                    │       ↓                    │
│                    │   PEP verifies:            │
│                    │   - Signature              │
│                    │   - Policy                 │
│                    │       ↓                    │
│                    │   Execute tool             │
└────────────────────┴────────────────────────────┘
\end{verbatim}}

\textbf{Security Challenge}: LLMs are untrusted for authorization decisions. Even with prompt engineering, adversarial prompts can cause unauthorized function calls. PEPs verify at function execution boundaries regardless of LLM decisions.

\subsection*{LangChain (Multi-Agent Orchestration)}

LangChain coordinates multi-step workflows across tools and LLMs. Multi-level integration addresses four layers:

{\tiny
\begin{verbatim}
LangChain Architecture:
┌─────────────────────────────────────────────────┐
│         LangChain Orchestrator                  │
│              (macawAgent)                       │
├──────────┬─────────┬──────────┬─────────────────┤
│ Agents   │ Tools   │  Chains  │    Memory       │
│          │         │          │                 │
│Coordinate│Individual│Sequential│Persistent state │
│workflow  │operations│compose   │across turns     │
│          │         │          │                 │
│Agent →   │Tool →   │Chain →   │Memory →         │
│macawAgent│ Tool    │ Policy   │AuthenticatedCntx│
│          │ (w/PEP) │ Intersect│ (hash chains)   │
│          │         │          │                 │
│Each agent│Each tool│Effective │Context integrity│
│has own   │verifies │policy =  │verified every   │
│identity  │CAW      │∩ of all  │state transition │
│and policy│before   │tool      │                 │
│          │execution│policies  │Prevents context │
│          │         │          │poisoning attacks│
└──────────┴─────────┴──────────┴─────────────────┘

Workflow Flow:
User: "Analyze documents and send summary"
  ↓
LangChain Agent (macawAgent)
  ↓
Step 1: Read docs
  → Tool invocation → CAW → PEP check → Execute
  ↓
Step 2: Analyze content
  → LLM invocation → CAW → PEP check → Execute
  ↓
Step 3: Send email
  → Tool invocation → CAW → PEP check → Execute

Policy Composition via Intersection:
If tool_1 policy allows {A, B, C}
And tool_2 policy allows {B, C, D}
And chain composes tool_1 → tool_2
Then effective policy = {B, C} (intersection)
Chain composition cannot relax tool policies.
\end{verbatim}}

\textbf{Security Challenge}: Orchestrators expose multiple control points beyond individual tools—chain composition logic, agent reasoning, memory operations. Securing only tools would leave composition and state management unprotected.

\subsection*{CrewAI (Role-Based Multi-Agent)}

CrewAI organizes agents into crews where each member has a role executing tasks collaboratively:

{\tiny
\begin{verbatim}
CrewAI Architecture:
┌─────────────────────────────────────────────────┐
│                   Crew                          │
│              (macawAgent)                       │
├──────────┬─────────┬──────────┬─────────────────┤
│ Members  │  Roles  │  Tasks   │  Collaboration  │
│          │         │          │                 │
│Agents    │Role-    │Work items│Inter-agent      │
│with      │based    │assigned  │communication    │
│specific  │perms    │to agents │                 │
│roles     │         │          │                 │
│          │         │          │                 │
│Member →  │Role →   │Task →    │Message →        │
│macawAgent│Attestat │ CAW      │ CAW with        │
│          │ ion     │          │ role check      │
│          │         │          │                 │
│Writer    │"writer" │Execute   │Writer cannot    │
│cannot    │attested │task only │invoke DB tools  │
│access DB │crypto-  │if role   │reserved for     │
│tools     │graphical│matches   │Researcher role  │
└──────────┴─────────┴──────────┴─────────────────┘

Role-Based Flow:
Crew: Researcher + Writer
  ↓
Task: "Write report on Q4 data"
  ↓
Researcher (macawAgent with "researcher" attestation):
  → Invokes database_query tool
  → PEP checks: signature valid? role="researcher"?
  → Policy: "database_query" requires attestation["researcher"]
  → Allowed: execute query
  ↓
Writer (macawAgent with "writer" attestation):
  → Receives data from Researcher
  → Attempts database_query tool
  → PEP checks: signature valid? role="writer"?
  → Policy: "database_query" requires attestation["researcher"]
  → Denied: Writer lacks required attestation
\end{verbatim}}

\textbf{Security Challenge}: Role-based permissions must be cryptographically enforced. A ``writer'' role should not access databases reserved for ``researchers,'' even if the LLM powering the writer attempts unauthorized invocations. Attestations bind roles cryptographically to agent identities.

\subsection*{AutoGen (Autonomous Code Generation)}

AutoGen enables autonomous code generation and execution. The security challenge is unrestricted code execution:

{\tiny
\begin{verbatim}
AutoGen Architecture:
┌─────────────────────────────────────────────────┐
│         AutoGen Conversation                    │
│              (macawAgent)                       │
├──────────┬─────────┬──────────┬─────────────────┤
│ Agents   │ Code    │  Exec    │   Validation    │
│          │  Gen    │          │                 │
│Conversat │Generate │Execute   │Verify code      │
│ional     │Python   │generated │before execution │
│agents    │code to  │code      │                 │
│          │solve    │          │                 │
│          │tasks    │          │                 │
│          │         │          │                 │
│Agent →   │LLM →    │Exec →    │Verifiers:       │
│macawAgent│ Tool    │ Tool     │- AST analysis   │
│          │         │ (w/PEP)  │- Allowlist check│
│          │         │          │- Sandbox verify │
│          │         │          │                 │
│Each agent│Code gen │Execution │Pre-invocation   │
│signs     │becomes  │request = │verifiers analyze│
│requests  │CAW with │CAW       │code content     │
│          │code     │          │before execution │
└──────────┴─────────┴──────────┴─────────────────┘

Code Execution Flow:
User: "Calculate Fibonacci(10)"
  ↓
AutoGen Agent generates code:
  def fib(n): return n if n<2 else fib(n-1)+fib(n-2)
  ↓
Execution request → CAW
  ↓
PEP invokes pre-execution verifiers:
  1. AST analysis: detect dangerous patterns
     (os.system, subprocess, eval, exec, import sys)
  2. Allowlist check: only math/basic operations
  3. Sandbox verification: no network/file access
  ↓
Verifiers approve → Execute in sandbox → Return result
  ↓
Malicious attempt: "os.system('rm -rf /')"
  ↓
AST analysis detects os.system → DENY
\end{verbatim}}

\textbf{Security Challenge}: Unrestricted code execution enables arbitrary operations including data exfiltration and privilege escalation. Pluggable verifiers analyze code through AST parsing before execution, enforcing allowlists and sandbox constraints.

\subsection*{LlamaIndex (RAG Pipelines)}

LlamaIndex provides RAG (Retrieval Augmented Generation) pipelines for document processing and question answering:

{\tiny
\begin{verbatim}
LlamaIndex Architecture:
┌─────────────────────────────────────────────────┐
│         LlamaIndex Pipeline                     │
│              (macawAgent)                       │
├──────────┬─────────┬──────────┬─────────────────┤
│  Index   │ Query   │ Retrieve │   Generate      │
│          │  Engine │          │                 │
│Document  │Process  │Vector    │LLM generates    │
│indexing  │queries  │search    │answer from      │
│          │         │          │retrieved docs   │
│          │         │          │                 │
│Index →   │Query →  │Retrieve →│Generate →       │
│macawAgent│ Tool    │ Tool     │ Tool            │
│          │ (w/PEP) │ (w/PEP)  │ (w/PEP)         │
│          │         │          │                 │
│Vector DB │Query    │Document  │LLM receives     │
│becomes   │validated│access    │only allowed     │
│macawAgent│against  │validated │documents        │
│          │policy   │against   │                 │
│          │         │policy    │                 │
└──────────┴─────────┴──────────┴─────────────────┘

RAG Flow with Security:
User: "What were Q4 earnings?"
  ↓
Query Engine (Tool with PEP)
  → CAW: query="Q4 earnings"
  → PEP checks:
     - Signature valid?
     - Policy allows query pattern?
     - Denied patterns: *password*, *credentials*
  → Approved: continue
  ↓
Retriever (Tool with PEP)
  → CAW: retrieve docs matching "Q4 earnings"
  → PEP checks:
     - Can access these documents?
     - Allowed files: *public*, *reports*
     - Denied files: *private*, *secret*
  → Approved: retrieve documents
  ↓
Generator (Tool with PEP)
  → CAW: generate answer from retrieved docs
  → PEP checks: output validation
  → Return answer to user
\end{verbatim}}

\textbf{Security Challenge}: RAG pipelines access untrusted document stores. Malicious documents could contain prompt injection or sensitive data. PEPs verify document access patterns and query constraints before retrieval.

\subsection*{Haystack (Document Processing Pipelines)}

Haystack provides document processing pipelines with nodes for retrieval, preprocessing, and generation:

{\tiny
\begin{verbatim}
Haystack Architecture:
┌─────────────────────────────────────────────────┐
│         Haystack Pipeline                       │
│              (macawAgent)                       │
├──────────┬─────────┬──────────┬─────────────────┤
│ Pipeline │  Nodes  │ Retrieve │   Process       │
│          │         │          │                 │
│Sequential│Pipeline │Document  │Text processing  │
│execution │steps    │retrieval │and generation   │
│of nodes  │         │          │                 │
│          │         │          │                 │
│Pipeline →│Node →   │Retrieve →│Process →        │
│macawAgent│ Tool    │ Tool     │ Tool            │
│          │ (w/PEP) │ (w/PEP)  │ (w/PEP)         │
│          │         │          │                 │
│Each node │Each node│Document  │Processing       │
│in chain  │verifies │access    │validated        │
│enforces  │CAW      │controlled│against policy   │
│policy    │         │by policy │                 │
└──────────┴─────────┴──────────┴─────────────────┘

Pipeline Flow:
User: "Process documents in ./reports/"
  ↓
Pipeline (macawAgent)
  ↓
Node 1: FileReader (Tool with PEP)
  → CAW: read("./reports/")
  → PEP checks:
     - Allowed paths: *reports*, *public*
     - Denied paths: *credentials*, *private*
  → Approved: read files
  ↓
Node 2: Preprocessor (Tool with PEP)
  → CAW: clean text, remove PII
  → PEP checks: PII detection verifier
  → Execute with redaction
  ↓
Node 3: Generator (Tool with PEP)
  → CAW: generate summary
  → PEP checks: output validation
  → Return result

Policy Composition Across Nodes:
If Node 1 allows {A, B, C}
And Node 2 allows {B, C, D}
And pipeline chains Node 1 → Node 2
Then effective policy = {B, C} (intersection)
\end{verbatim}}

\textbf{Security Challenge}: Document processing pipelines access file systems and external data sources. Nodes must enforce path constraints and content validation to prevent unauthorized access or data exfiltration.
